\subsection{Konzept der Persona-Entwicklung}

Die Analyse der durchgeführten Interviews, der Fragebögern und Sekundärdaten bildete die Grundlage für die Erstellung von Personas. Personas sind detaillierte Archetypen von Zielnutzer:innen, die auf Forschung basieren. Sie helfen dem Team, die Perspektive der Zielgruppe zu verstehen und Lösungen nutzer-zentriert zu entwickeln.

Der Ablauf der Persona-Entwicklung war folgend:
Der Ablauf der Persona-Entwicklung folgte einer klaren Sequenz: Zuerst wurden die Interviews aus der Discovery-Phase ausgewertet und die Interviewpersonen entlang kontinuierlicher und diskreter Variablen gemappt. Darauf aufbauend erfolgte ein Clustering der Teilnehmenden, aus dem die Personas entwickelt wurden. Anschliessend wählten wir eine Leit-Persona für die vertiefte Analyse, erstellten für diese eine Customer Journey, verdichteten die Erkenntnisse in einem Point of View und leiteten daraus die How-might-we-Fragen ab.

\subsection{Persona: Kurt Kreuz und Quer}

Verdichtet aus Interviews: Reini (66), Hans Jörg (78), Peter (66)

\textit{``Standard bei mir ist die \gls{oev}. Zug, Tram, Bus, das macht doch einfach Sinn. Das Auto kommt zum Einsatz, wenn ich muss, schwere Einkäufe, spät in der Nacht ohne Anschluss. Ich bin in Vereinen tätig und gesellig, das geht im \gls{oev} halt auch am besten.''}

\subsubsection{Demografische Merkmale}

\end{multicols}

\begin{center}
\begin{tabular}{ll}
Alter & 65 bis 80 Jahre (frisch pensioniert) \\
Wohnort & Seegemeinde / städtische Agglo Zürich \\
Beruf (i.R.) & Technisch / handwerklich (Elektriker, Mechaniker, Ing.) \\
Mobilität & \gls{oev} (Freizeit \& Alltag) · Auto (Ausnahme) \\
\gls{oev}-Abo / Digital & Halbtax oder \gls{ga} · \gls{sbb}-App, \gls{zvv}-App, MeteoSwiss, WebCams
\end{tabular}
\end{center}

\begin{multicols}{2}

\subsubsection{Bedürfnisse \& Ziele}

\textbf{Bewegung \& Gesundheit:} Körperlich und geistig fit bleiben. Neues entdecken, Neues lernen, Gehirnjogging als bewusste Strategie gegen Stillstand.

\textbf{Nachhaltig handeln:} Umweltschutz ist wichtig. Möchte als Vorbild vorangehen und zeigen, dass man auch ohne Auto auskommen kann.

\textbf{Geselligkeit \& Sinnhaftigkeit:} Aktiv in Vereinen, gibt Wissen an Jüngere weiter. Enttäuscht über Ausschluss wegen Alter in manchen Kontexten.

\textbf{Entdecken \& Erleben:} Orte, Natur, Kultur, Menschen. Die Fahrt ist Erlebnis, nicht nur Transport. Wissen wird auf dem Weg gesammelt.

\subsubsection{Einstellungen}

\textbf{\gls{oev} als natürliches Default:} \gls{oev} ist erste Wahl für alle Wege. Auto kommt nur bei echtem Nutzen: schwere Einkäufe für den Garten, späte Nacht ohne Anschluss.

\textbf{Kritisch bei Komplexität:} Grundsätzlich \gls{oev}-affin, ärgert sich aber über komplizierte Zonensysteme und unlogische Ticketregeln.

\subsubsection{Verhalten}

\textbf{Tracking \& Dokumentieren:} Misst zurückgelegte Kilometer, Bergsteigungen, besuchte Orte. Sammeln und Messen ist Gewohnheit, motiviert und hält den Überblick.

\textbf{Recherche \& Eigeninitiative:} Bei unklaren \gls{oev}-Systemen wird selbst recherchiert; Ausflüge werden mit Höhenmetern und Webcams geplant, das macht Spass und hält geistig fit.

\textbf{Verantwortung \& Vorbild:} Bringt sich in Vereine ein, plant Aktivitäten mit, zeigt anderen wie es geht. Möchte Vorbild sein, auch beim Mobilitätsverhalten. ``Reini, machst du das?''

\subsection{Customer Journey: \gls{oev}-Ausflug mit der Gruppe}

Eine Customer Journey beschreibt den gesamten Weg, den eine Person durchläuft, wenn sie ein Produkt oder eine Dienstleistung nutzt, von dem Moment, wo die Idee entsteht, bis zur Nachbetrachtung des Erlebnisses. Der Begriff kommt aus dem Design Thinking und dem Marketing. Die Idee dahinter ist simpel: Man versetzt sich konsequent in die Perspektive des Nutzers, statt aus der Innenperspektive des Unternehmens zu denken.

Wir wählten die typische \gls{oev}-Nutzung der Persona ``Kurt'' ab: vom Gedanken eine Reise unternehmen zu wollen, bis zum Abschluss der Reise mit der Reflexion.

\subsubsection{Reise-Phasen und Touchpoints}

Die Journey beginnt mit der \textbf{Inspiration} aus Eigeninitiative, häufig ausgelöst durch Medien, Vereinskontexte oder Gespräche im Freundeskreis. In der anschliessenden \textbf{Koordination} entsteht bereits ein positiver sozialer Effekt, weil die Organisation gemeinsamer Aktivitäten als sinnstiftend erlebt wird. In der \textbf{Detailplanung} zeigt sich ein ambivalentes Bild: Digitale Recherche, Reservationen und Routenabgleich erzeugen Entdeckerfreude, gleichzeitig entstehen Friktionen bei Zusatzangeboten und der Integration von Bergbahn- oder Freizeitleistungen.

Besonders kritisch ist der Schritt des \textbf{Ticketkaufs}, da die Tariflogik für Gruppen als komplex und fehleranfällig wahrgenommen wird. Der \textbf{Weg zum Bahnhof} wird dagegen eher neutral bis positiv erlebt, weil Bewegung und der Verzicht auf Parkplatzsuche Vorteile bieten. Während der \textbf{Hinfahrt} überwiegen Landschaftserlebnis, Gespräche und Entspannung, wobei uneindeutige Umstiege punktuell Stress erzeugen können. Das eigentliche \textbf{Erlebnis am Zielort} ist emotional der stärkste Moment der Journey und bestätigt den Nutzen der gemeinsamen \gls{oev}-Reise.

Die \textbf{Rückreise} wird meist routinierter und ruhiger wahrgenommen, bleibt aber störungsanfällig bei Verspätungen und engen Anschlüssen. In der abschliessenden \textbf{Reflexion} dominiert Stolz auf die Organisation des Ausflugs, was als wichtiger Hebel für Wiederholung und soziale Weiterempfehlung verstanden werden kann.

\subsection{Point of View}

Basierend auf der Persona und der Customer Journey wird das Problemverständnis in einem Point of View verdichtet:

\begin{quote}
\textit{``Ich kenne den \gls{oev} in- und auswendig und merke, dass viele in meinem Umfeld gar nicht wissen, wie einfach das ist. Es macht mir Freude, andere mitzunehmen, die Route zu planen und Verantwortung zu übernehmen, damit sie einfach geniessen können. Ich möchte der Grund sein, warum jemand das nächste Mal selbst zum \gls{oev} greift.''}
\end{quote}

Zentrales Insight: Kurt ist nicht der Zielkunde, Kurt nutzt den \gls{oev} bereits. Kurt ist das Aktivierungspotenzial. Der wirkliche Zielkunde ist der Freund oder die Freundin, die Kurt überredet, mal einen Ausflug mit dem \gls{oev} zu machen. Dieses Segment zu aktivieren, durch Botschafter wie Kurt, könnte eine kritische Massnahme sein.

\subsection{How-might-we-Fragen (Impulsfragen)}

Aus der Analyse entstanden mehrere How-might-we-Fragen, die den Innovationsraum für Lösungsideen aufspannen:
Im Zentrum standen dabei sechs Leitfragen: Wie können versierte \gls{oev}-Nutzer Ausflüge für Gruppen einfacher planen und organisieren? Wie lassen sich Nicht-\gls{oev}-Nutzende so einbinden, dass sie ohne Planungsaufwand teilnehmen können? Wie kann der erste gemeinsame \gls{oev}-Ausflug so gestaltet werden, dass ein bleibender positiver Eindruck entsteht? Wie lassen sich erfahrene Nutzer:innen motivieren, als Botschafter:innen in ihrem Umfeld aktiv zu werden? Wie kann der Ankommensmoment ohne Parkplatzsuche als emotionales Plus sichtbar werden? Und wie vereinfachen wir Ticketing und Verbindungswahl für ganze Gruppen in einem Schritt?

Diese Fragen bilden den Ausgangspunkt für die Develop und Deliver Phase.

\subsection{Weitere Personas im Spektrum}

Neben Kurt Kreuz und Quer identifizierten wir zwei weitere zentrale Personas im Spektrum:

\subsubsection{Persona 2: Maria Mobil (65 bis 75 Jahre, ländlicher Raum)}

\textit{``Ich bin gerne unterwegs, aber mit dem Auto bin ich einfach flexibler. Der Bus kommt nur stündlich, und wenn ich irgendwo um 16 Uhr sein muss, schaffe ich es oft nicht.''}

\textbf{Charakteristiken:} Aktive, kaufkräftige Seniorin im ländlichen Raum. Das \gls{oev}-Angebot ist nicht attraktiv genug, lange Wartezeiten, keine Direktverbindungen. Hürde ist nicht psychologisch, sondern systemisch. Sie würde \gls{oev} nutzen, wenn das Angebot besser wäre.

\textbf{Implikation für Massnahmen:} Bei Maria braucht es nicht primär Bewusstseinswandel, sondern bessere Service-Angebote oder zielgerichtete Informationen über beste Verbindungen.

\subsubsection{Persona 3: Heinz Hüfthoch (78 Jahre, betagter Senior)}

\textit{``Seit ich nicht mehr Auto fahre, gehe ich weniger raus. Mit dem Bus bin ich unsicher, ich weiss nicht, wo ich umsteigen soll.''}

\textbf{Charakteristiken:} Betagter Senior, hat kürzlich Führerschein aufgegeben. Mobilitätsverlust ist emotionaler Schock. Hauptbarriere: Orientierungsunsicherheit, Angst vor Fehlern, Sicherheitsbedenken.

\textbf{Implikation für Massnahmen:} Heinz braucht intensive persönliche Begleitung, niederschwellige Unterstützung, und möglicherweise ein ``Buddy-System'' mit jemandem erfahrenerem. Der kritische Moment der Führerscheinabgabe ist für ihn besonders wichtig.

\subsection{Persona-Matrix und Priorisierung}

Die drei Personas spannen ein Spektrum auf:

\end{multicols}

\begin{center}
\begin{tabular}{l|l|l|l}
& Kurt (aktiv, urban, \gls{oev}-affin) & Maria (aktiv, ländlich, \gls{oev}-skeptisch) & Heinz (betagter, unsicher, Transition) \\
\hline
Alter & 65 bis 75 & 65 bis 75 & 75 bis 85 \\
Wohnort & urban & ländlich & urban/suburban \\
Status & \gls{oev}-Nutzer & Auto-Nutzer & Transition \\
Aktivierungshebel & Botschafter & Service/Info & Begleitung \\
\end{tabular}
\end{center}

\begin{multicols}{2}

Priorität für Massnahmen-Design:
Für das Massnahmen-Design priorisieren wir drei Hebel in abgestufter Reihenfolge: Erstens Kurt als Botschafter mit Multiplikatoreffekt, zweitens Maria mit Fokus auf Service-Optimierung und alltagsnahe Best-Practice-Informationen und drittens Heinz mit gezielter Begleitung im Übergang nach der Führerscheinabgabe.
